\begin{proposition}[DP30: a dual direction with a unit $p$-th derivative]\label{DP30}
Let $A$ be a complete mixed-characteristic discrete valuation ring with perfect residue field $k$. Let $G/A$ be a one-dimensional strict formal $\mathcal O$-module, where $\mathcal O\subset A$ is the ring of integers of the unramified degree-$h$ extension of $\Q_p$, and suppose its special fibre has height exactly $h\ge2$. Let $H=G[p^\infty]^D$ be its actual Cartier dual and $\widehat H$ its connected formal group. Then $\Lie(H_k)$ has a basis $e_0,\ldots,e_{h-2}$ with
\[
 e_i^{[p]}=e_{i+1}\ (i<h-2),\qquad e_{h-2}^{[p]}=0.
\]
Thus the restricted $p$-operation has kernel dimension one and image dimension $h-2$. If $h>2$, there are a primitive integral invariant derivation $D$ on $\widehat H$ and an integral formal coordinate $f$ such that $D^pf(0)\in A^\times$.
\end{proposition}
\begin{proof}
Teichm\"uller averaging as in DP19 gives an integral coordinate $x$ with $[\zeta](x)=\zeta x$ for every $\zeta^{p^h-1}=1$. The group law then satisfies $F(\zeta X,\zeta Y)=\zeta F(X,Y)$, so its nonzero monomials have total degree congruent to one modulo $q-1$, where $q=p^h$. In particular $F=X+Y+$ terms of degree at least $q$, integrally and after Frobenius twisting.

For the next calculation work over $k$. Relative Frobenius $F_G:G\to G^{(p)}$ and Verschiebung $V_G:G^{(p)}\to G$ compose in either order to multiplication by $p$; their kernel orders are $p$ and $p^{h-1}$. These facts and the duality of $F,V$ are Tate \cite[\S2.2, diagram (3), pp.~163--164; \S2.3, Proposition 3, p.~166]{Tate67}. Restrict $V_G$ to $G^{(p)}[p]$. Its image has order $p$ and lies in $G[F_G]$, so it is exactly $G[F_G]$. Exact finite Cartier duality \cite[\S1.2--1.3, pp.~159--160]{Tate67} gives
\[
 H[F_H]^D\simeq G[p]/G[F_G].
\]
Height $h$ implies $[p](x)=u(x)x^q$ with $u(0)\ne0$, so $G[p]=G[F_G^h]$ as subgroup schemes, by equality of their defining ideals. If $r=p^{h-1}$, relative Frobenius induces
\[
 G[p]/G[F_G]\simeq G^{(p)}[F^{h-1}].
\]
On coordinates this is $k[z]/(z^r)\hookrightarrow k[x]/(x^q)$, $z\mapsto x^p$. The target is free of rank $p$ with basis $1,x,\ldots,x^{p-1}$, proving faithful flatness and the asserted quotient. In $k[X,Y]/(X^r,Y^r)$ every surviving monomial has degree at most $2r-2<q$. The nonlinear part of the twisted group law vanishes. Hence
\begin{equation}\label{eq:dual-F-kernel}
 H[F_H]^D\simeq\alpha_{p^{h-1}}.
\end{equation}

Put $P=H[F_H]$. Its distribution Hopf algebra is, by duality,
\[
 \Dist(P)=k[t]/(t^r),\qquad\Delta(t)=t\otimes1+1\otimes t.
\]
The Lie algebra is its primitive subspace: the augmentation Leibniz rule for a tangent functional is precisely the primitive-element identity in the dual Hopf algebra. This subspace has basis $t,t^p,\ldots,t^{p^{h-2}}$. Indeed these powers are primitive, while for any other exponent $1\le n<r$ a nonzero base-$p$ digit gives an interior binomial coefficient $\binom n{p^j}\not\equiv0\pmod p$. The resulting mixed tensor term survives; different total degrees cannot cancel it. Constants are not primitive.

The restricted operation is the $p$-th power in this algebra: convolution powers of a tangent functional correspond to powers of its invariant derivation. It shifts the displayed basis and kills its last member. Since $dF_H=0$, the inclusion $P\hookrightarrow H$ identifies their Lie algebras and their restricted operations. Perfectness of $k$ gives the stated image and kernel dimensions.

The connected formal group exists by Tate \cite[\S2.2, Proposition 1, p.~162 and (4), p.~164]{Tate67}; its tangent lattice is free of rank $h-1$, since dual dimensions sum to height by Proposition 3. For $h>2$, lift a tangent vector $v$ with $v^{[p]}\ne0$ to this lattice and let $D$ be its invariant derivation. The lift is primitive. Some integral coordinate $f$ has differential detecting $v^{[p]}$. Reduction commutes with $D$ and its iterates, so
\[
 \overline{D^pf(0)}=df_0(v^{[p]})\ne0.
\]
Thus $D^pf(0)$ is a unit. No assertion that the full $H$ is connected is needed for this last argument.
\end{proof}
